% Exported from Math++. Compile with pdflatex.
\documentclass[11pt,a4paper]{article}
\usepackage[T1]{fontenc}
\usepackage[utf8]{inputenc}
\usepackage{lmodern}
\usepackage{amsmath,amssymb,amsthm,mathtools}
\usepackage{graphicx,booktabs,xcolor}
\usepackage[hidelinks]{hyperref}
\usepackage[margin=25mm]{geometry}
\theoremstyle{plain}
\newtheorem{theorem}{Theorem}
\newtheorem{lemma}[theorem]{Lemma}
\newtheorem{corollary}[theorem]{Corollary}
\newtheorem{proposition}[theorem]{Proposition}
\theoremstyle{definition}
\newtheorem{definition}[theorem]{Definition}
\newtheorem{example}[theorem]{Example}
\theoremstyle{remark}
\newtheorem*{remark}{Remark}
\setlength{\parindent}{0pt}\setlength{\parskip}{0.6em}
\title{Algebra}
\date{}
\begin{document}
\maketitle
Expanding, factoring and solving work the way you would do them on paper.

\section*{Expand and factor}

\begin{equation*}
  \operatorname{expand}((x + 1)^3) = x^{3} + 3 x^{2} + 3 x + 1
\end{equation*}
\begin{equation*}
  \operatorname{factor}(x^2 - 5 x + 6) = \left(x - 3\right) \left(x - 2\right)
\end{equation*}
\begin{equation*}
  \operatorname{factor}(x^4 - 1) = \left(x - 1\right) \left(x + 1\right) \left(x^{2} + 1\right)
\end{equation*}
\section*{Solve}

The roots of a quadratic:

\begin{equation*}
  \operatorname{solve}(x^2 - 5 x + 6 = 0, x) = \left\{2, 3\right\}
\end{equation*}
Give a name to a value and use it later:

\begin{equation*}
  a \coloneqq  3
\end{equation*}
\begin{equation*}
  b \coloneqq  4
\end{equation*}
\begin{equation*}
  \sqrt{a^2 + b^2} = 5
\end{equation*}
\begin{equation*}
  \frac{1}{a} + \frac{1}{b} = \frac{7}{12}
\end{equation*}
\begin{equation*}
  \sqrt{2} \cdot \sqrt{8} = 4
\end{equation*}
\end{document}
